同一份数据，两个人用同一个公式算方差，结果差了六位有效数字。

公式没错，两人写的都对：一个用 \(\E[X^2]-(\E X)^2\)，另一个逐点累加偏差。数学上它们恒等，浮点里不是——前者在均值远大于标准差时，两个相近的大数相减，有效位数被一次抵消吃掉大半。这不是谁写错了代码，而是\emph{同一个数学表达式有无穷多种等价写法，它们的数值行为却完全不同}。

这句话是整个数值分析的入口。前面几部分关心公式什么时候成立；这一部分关心公式成立之后，算出来的那个数能不能信。两者之间隔着三样东西：机器只有有限位数、问题本身可能对扰动敏感、算法自己还会添乱。

三者必须分开记账，否则力气会花错地方。

\begin{center}
\begin{tikzpicture}[x=1cm,y=1cm,font=\small,>=stealth]
  \draw[black!45,fill=blue!8] (0,2.1) rectangle (4.0,3.3);
  \node[align=center] at (2.0,2.7) {机器只有有限位数};
  \node[align=center,font=\footnotesize,black!65] at (2.0,1.6) {相对误差 \(\approx\varepsilon\)\\消不掉，只能不放大};
  \draw[black!45,fill=orange!9] (5.0,2.1) rectangle (9.0,3.3);
  \node[align=center] at (7.0,2.7) {问题有多敏感};
  \node[align=center,font=\footnotesize,black!65] at (7.0,1.6) {条件数 \(\kappa\)\\与用什么算法无关};
  \draw[black!45,fill=green!8] (10.0,2.1) rectangle (14.0,3.3);
  \node[align=center] at (12.0,2.7) {算法添了多少乱};
  \node[align=center,font=\footnotesize,black!65] at (12.0,1.6) {后向误差\\只有这一项归你管};
  \draw[black!40] (0,0.55) -- (14.0,0.55);
  \node[align=center] at (7.0,0.0) {最终误差 \(\;\approx\;\) 后向误差 \(\times\) 条件数};
  \draw[->,black!50] (2.0,2.05) -- (2.0,0.62);
  \draw[->,black!50] (7.0,2.05) -- (7.0,0.62);
  \draw[->,black!50] (12.0,2.05) -- (12.0,0.62);
\end{tikzpicture}

\emph{病态问题上换再好的算法也救不回精度，良态问题上写坏的实现照样出错。先判断自己在哪一格，再决定动手改什么。}
\end{center}

\subsection{四条贯穿始终的判断}

\textbf{有意义的是相对误差，不是绝对误差。}浮点数在数轴上不是均匀分布的，而是按指数分段、每段内等距，所以``差了 \(0.001\)''这句话脱离量级就没有含义。机器精度刻画的正是相对误差的下限，后面一切误差分析都用这把尺子。

\textbf{条件数属于问题，稳定性属于算法。}这两个概念被混淆得最多。条件数说的是输入扰动被放大多少倍，与你用什么方法无关；稳定性说的是方法自身引入多少额外误差。训练不收敛时，先分清是数据共线让问题本身病态，还是实现里该用 log-sum-exp 却直接算了指数——两种情况该动的东西完全不同。

\textbf{残差小不等于解准。}方程几乎被满足，与未知量几乎正确，中间隔着一个条件数。这条判断在最小二乘（正规方程把条件数平方）、迭代求解器（残差降到 \(10^{-10}\) 而解只剩几位有效数字）、优化的收敛判据里反复出现。检查工作时不要只看残差。

\textbf{结构就是可以兑现的计算资源。}对称正定省掉一半运算与全部选主元，稀疏把直接法整个换成迭代法，光滑性让谱方法的误差按指数而非多项式下降。每次利用结构都是在用一个假设换代价，所以也要问：假设不成立时会以什么形式暴露出来。谱方法遇到间断、Krylov 遇到分散的谱、插值遇到噪声数据，都是这个问题的不同版本。

\subsection{九章怎么安排}

\hyperref[ch:061]{``浮点数与误差：计算机为什么不等于实数系统''}是地基，必须先读。它建立的三个概念——相对误差、条件数、后向稳定——是后面八章共用的语言。

\hyperref[ch:065]{``数值线性代数：结构、稳定性与规模决定算法''}使用频率最高，做数据工作的读者几乎每天在间接使用它的结论。\hyperref[ch:067]{``机器学习中的稳定计算''}最贴近日常工程，篇幅短，可以提前读——它把前面的原则落成三个具体改写：换到对数域、换到计算图、换成解方程。

\hyperref[ch:062]{``非线性方程求根：可靠性、速度与信息成本''}与\hyperref[ch:063]{``插值与逼近：通过数据点，不等于理解点之间''}给出两组反复出现的权衡：慢而可靠对快而可能失手、穿过每个点对整体接近。后者的 Runge 现象是全书关于``容量更大未必更好''的最早例子。

\hyperref[ch:064]{``数值微分与积分：同一份采样，为何难度相反''}用同一份采样点做两件事，一件天生放大误差、一件天生平均误差。\hyperref[ch:066]{``常微分方程数值方法：沿局部斜率走完整条轨迹''}讲清误差累积与稳定性限制，其中刚性方程一篇解释了``步长被最快模态绑架''，与优化里学习率被最敏感方向绑架是同一个结构。

\hyperref[ch:068]{``Fourier 与谱方法：把函数换到频率坐标中''}与\hyperref[ch:069]{``复分析与留数：奇点决定函数的深层结构''}适合连读：前者说明换到频率坐标买到了什么、付出了什么，后者给出系数衰减快慢背后的原因——函数在复平面上离实轴多远才出现奇点。Runge 现象在那里得到最终解释。

读这一部分时保持一个习惯：拿到一个数值结果，先问它的误差来自哪一笔账。这个追问不需要额外知识，却把``跑出来了''和``可以拿去用''分成了两件事。
